% 03-related.tex
\section{Related Work}
\label{sec:related}

\paragraph{Component-level analysis of sidecars}
MeshInsight~\cite{zhu2023meshinsight} is the work closest in method to ours.
It profiles the Istio sidecar (Istio 1.13, Envoy 1.21) on Online Boutique and
Hotel Reservation and attributes its overhead to seven components, finding that
protocol parsing dominates. It reports latency increases of 27--269\% and CPU
increases of 42--163\%, depending on configuration, and proposes a linear
compositional model of overhead. The authors note that the model does not
predict tail latency, and that extending the analysis to other meshes is future
work. We extend the analysis to ztunnel, the Ambient waypoint,
\texttt{linkerd2-proxy}, and Cilium's eBPF datapath, use open-loop load, add
energy, and model the tail through a queueing term.

\paragraph{Comparisons of sidecar and sidecar-less meshes}
Bremler-Barr et al.~\cite{bremlerbarr2024mtls} compare Istio sidecar, Istio
Ambient (ztunnel only), Linkerd, and Cilium, focusing on the cost of mTLS. Their
workload is a synthetic server with a fixed 200\,ms delay driven by Fortio, and
they do not exercise L7 processing, waypoints, or realistic call graphs. They
also observe that Cilium does not encrypt traffic between pods on the same node,
a caveat that applies to our setup as well: under our placement, 2.8 of the
3.8 internal calls per Hotel Reservation request, and 0.6 of the 11 per Online
Boutique request, stay on one node and are therefore not encrypted by Cilium's
WireGuard. Our study uses multi-service
applications, includes the waypoint, and reports energy and a predictive model.

\paragraph{Reducing mesh overhead}
A separate line of work proposes new data-plane designs rather than measuring
existing ones. Meshlib~\cite{mostafaei2026meshlib} builds on Cilium and enforces
L7 policy inside the application process through an optional library, falling
back to Cilium's per-node Envoy for unmodified services; on TrainTicket, a
single node and up to 200\,req/s, it reports lower latency than Istio, Linkerd,
and unmodified Cilium, and attributes most of the cost of L7 proxying to
protocol parsing. Cvetkovi\'c and Klimovic~\cite{cvetkovic2025lightweight}
characterize Istio and Linkerd sidecars on Knative and report 38.3$\times$
higher p95 execution time with Istio policies and up to 3.9$\times$ higher p95
CPU per worker; they could not evaluate Ambient, which Knative does not
support. Neither evaluates Istio Ambient or energy. Our measurements are complementary: they quantify the cost these
designs aim to remove, on the architectures in production today.

\paragraph{Microservice benchmarks and load generation}
We use Online Boutique~\cite{onlineboutique} and Hotel Reservation from
DeathStarBench~\cite{gan2019deathstarbench}, the same applications as
MeshInsight, so that results can be compared directly. We generate open-loop
load with k6~\cite{k6} at a constant arrival rate, avoiding the
coordinated-omission problem of closed-loop generators, which underestimates
tail latency~\cite{tene2015latency}. Energy is measured with Intel
RAPL~\cite{khan2018rapl}.
Container-level energy estimators for Kubernetes such as Kepler have been
found inaccurate~\cite{pijnacker2025kepler}, and per-task RAPL attribution on
shared nodes must fall back to approximations from CPU
utilization~\cite{thamm2025rapl}; we therefore report package-level energy per
request for the whole host rather than per container.
